% answers_short.tex —— 1977 年·江苏南通卷 · 简短答案速查
% 由 tools/make_short_answers.py 自 answers.tex 的 \gksolhead/\gkanswer 提取，
% 并与 1977/1977_数学真题_江苏南通_答案要点.md 逐题比对后生成，勿手改（改源 MD 后重跑）。

\begin{gknotice}
\textbf{使用说明}：本节只给最终答案便于快速核对，完整推导见解析版。答案为编写者独立推导，\textbf{非官方标准答案}。
\end{gknotice}

\gkpart{题目}
\begin{itemize}[leftmargin=2.6em]
  \item[\textbf{1.}] (1) $a>0$ 时 $2a>a$；$a=0$ 时 $2a=a$；$a<0$ 时 $2a<a$；(2) $x\ge 0$
  \item[\textbf{2.}] (1) $x\ne 2$；(2) $x\ge\dfrac{1}{2}$
  \item[\textbf{3.}] $3\dfrac{7}{36}$（即 $\dfrac{115}{36}$）。
  \item[\textbf{4.}] $k=\pm 2$；$k=2$ 时两相等实数根为 $x_1=x_2=2$，$k=-2$ 时两相等实数根为 $x_1=x_2=-2$。
  \item[\textbf{5.}] $\cos\alpha=-\dfrac{4}{5}$，$\operatorname{tg}\alpha=-\dfrac{3}{4}$，$\operatorname{ctg}\alpha=-\dfrac{4}{3}$（即 $-1\dfrac{1}{3}$）。
  \item[\textbf{6.}] 顶点坐标 $(1,3)$，对称轴方程 $x=1$，函数 $y$ 有最大值 $3$。
  \item[\textbf{7.}] $x=8$；$x=2$ 是增根，应舍去。
  \item[\textbf{8.}] $3366$。
  \item[\textbf{9.}] 侧棱长为 $10\ \mathrm{cm}$。
  \item[\textbf{10.}] $F_1=\dfrac{P\sin\beta}{\sin(\alpha+\beta)}$，$F_2=\dfrac{P\sin\alpha}{\sin(\alpha+\beta)}$。
\end{itemize}
